Our conclusions for arrival sequence do not apply when evidence is hard or
updating is Bayesian. What an impression fixes is therefore the modelling
commitment on which everything here turns, and it admits two readings of the
same cue. On the first, the credential resets the evaluator's credence in
competence to a \emph{level}---the panelist comes to feel the candidate
competent to some degree---and the cue is read off as a delivered marginal on
that attribute's partition. On the second, the credential carries a
\emph{Bayes factor}---a likelihood ratio recording how many more times the
credential is to be expected under competence than in its absence---and the cue
is read off as multiplicative evidence to be applied against whatever belief it
meets. The two readings are not interchangeable. A delivered marginal is a
level an evaluator can plausibly report; a Bayes factor requires a counterfactual
likelihood---the probability of this credential were the candidate \emph{not}
competent---that impression-formation on soft input does not furnish. For input
of that kind strict conditionalisation is unavailable in any case, since there
is no certifiable proposition to condition on. We therefore take impressions to
arrive as revised credences on a cue's partition, and the phenomenology of the
example supports the marginal reading over the Bayes-factor one: the panel
feels a level, it does not compute a ratio.

The two readings are precisely the two objects the paper already compares. Read
as delivered marginals, the cues are composed by Jeffrey conditioning, which
holds conditionals fixed and---as a mechanism revising the prior to the
delivered marginal---is the unique coherent such revision and the minimal change
of the prior consistent with it, in the $I$-divergence sense
\citep{DiaconisZabell1982}\footnote{See \citet{Field1978} for how Jeffrey
updates are asymmetric in the delivered-credence parametrisation and
\citet{Wagner2002} for how they commute when their inputs are held fixed in
Bayes-factor form.}. Read as Bayes factors, the same two cues compose to the
sequence-free benchmark $\PB$ of Section~\ref{sec:jeffrey}, which commutes by
construction \citep{Wagner2002}. The order effect is thus not a property of soft
evidence as such but the gap between these two readings of identical cues: it is
present under the marginal reading and absent under the Bayes-factor one, and
the choice between them is what the phenomenon turns on. That order effects in
impression formation are a documented regularity \citep{Asch1946,HogarthEinhorn1992}
while Bayes-factor updating predicts none \citep{Wagner2002} is itself a reason
to prefer the marginal reading: Jeffrey conditioning is the minimal coherent
mechanism under which the established behavioural phenomenon can arise at all.
